\section{Primary-source audit: exact statements and boundaries}
This is a targeted audit, not a claim that all relevant literature has
been exhaustively searched. The following locators were checked in the
primary texts or their attached source files during this run.

\paragraph{Kadish--Landsberg.}
\emph{Padded polynomials, their cousins, and geometric complexity
theory}, Communications in Algebra42 (2014),2171--2180,
\href{https://arxiv.org/abs/1204.4693}{arXiv:1204.4693},
Theorems1.2,1.3,1.7 and Proposition1.8.
The necessary padding restrictions and the sufficient branch
$p_1\ge dn-m$ are retained. S-009 gives a complete counterexample
to the other sufficient branch $p_1\ge d(n-1)$ in Theorem1.3 of
the inspected arXiv version1, printed p.2. The published full text
was inaccessible, so this correction is expressly version-scoped.
Theorem1.7 describes a generalized Foulkes-map kernel, while
Proposition1.8 describes the relevant normalization under its hypotheses.
Remark1.4 explicitly allows an ideal to contain only some copies of a
Schur module. None of these statements computes the two determinant
evaluation ranks required in the main question. The main article was
read online because the attached KL directory contains definitions and
bibliography but not its main source.

\paragraph{B\"urgisser--Ikenmeyer--Panova and Ikenmeyer--Panova.}
\emph{No occurrence obstructions in geometric complexity theory},
\href{https://arxiv.org/abs/1604.06431}{arXiv:1604.06431v3},
Theorem1.4, excludes occurrence obstructions in its range $n\ge m^{25}$;
section1(c) explicitly retains multiplicity obstructions as a possibility.
Theorem2.1 and Theorem4.9 specify the length and padding restrictions.
The paper's displayed permanent uses internal padding; the separate
variable convention must be handled as in S-004 rather than silently
importing its smaller support cap.
\emph{Rectangular Kronecker coefficients and plethysms in geometric
complexity theory}, Advances in Mathematics319 (2017),40--66,
\href{https://arxiv.org/abs/1512.03798}{arXiv:1512.03798v2},
Theorem1.6 treats ordinary rectangular Kronecker vanishing for
$n>3m^4$; Corollary1.9 gives its degree lower bound. These are not
equalities for actual determinant orbit-closure multiplicities.

\paragraph{D\"orfler--Ikenmeyer--Panova.}
\emph{On geometric complexity theory: multiplicity obstructions are
stronger than occurrence obstructions},
\href{https://arxiv.org/abs/1901.04576}{arXiv:1901.04576v1},
Theorem2.3, compares Chow varieties with bounded-Waring-rank varieties.
Its finite case at degree7 and weight $(34,6,2)$ gives multiplicities
7 and8. Proposition3.2 and section6 contain the explicit finite
calculation; Theorem3.5 supplies its infinite plethysm identity.
These are genuine precedents for full multiplicity comparisons, but
the varieties are different from the padded-permanent/determinant pair.

\paragraph{Ikenmeyer--Kandasamy.}
\emph{Implementing geometric complexity theory: on the separation of
orbit closures via symmetries},
\href{https://arxiv.org/abs/1911.03990}{arXiv:1911.03990v1},
Proposition4.1, Theorems4.2 and4.3, and section5.
Theorem4.2 transports specified power-sum orbit multiplicities to its
closure after a rectangular shift. Theorem4.3 gives the power-sum versus
product-of-variables multiplicity obstruction, with section5 treating
the stronger distinctions from occurrence tests. This transport theorem
does not directly identify the determinant orbit-closure restriction
image needed here.

\paragraph{Kumar and B\"urgisser--Ikenmeyer--H\"uttenhain.}
Kumar, \emph{A study of the representations supported by the orbit
closure of the determinant}, Compositio Mathematica151 (2015),292--312,
\href{https://kumar.math.unc.edu/papers/kumar67.pdf}{author PDF},
Theorem6.1 and Corollaries6.2,6.4,6.5, with Remark6.6(b).
The occurrence results include even-size/Alon--Tarsi hypotheses where
stated. B\"urgisser--Ikenmeyer--H\"uttenhain,
\emph{Permanent versus determinant: not via saturations},
\href{https://arxiv.org/abs/1501.05528}{arXiv:1501.05528},
Theorem1, treats saturation for partitions of length at most$n$ and
size divisible by$n$; Theorem2 states the conditional strengthening.
These statements do not supply A-030's length-$n^2$ explicit Hessian
equation or A-034's exact ray algebra. This is a scope comparison,
not an exhaustive novelty certificate for the inherited claims.

\paragraph{Landsberg's book.}
The inspected text is the author's 2016 prepublication version of
\emph{Geometry and Complexity Theory},
\href{https://people.tamu.edu/~jml/simonsclass.pdf}{author PDF};
its pagination differs from the published2017 edition.
Propositions8.4.1.1 and8.4.2.1--8.4.2.4 give subspace and padding
restrictions and explicitly use the independent-variable cap$m^2+1$.
Theorem8.5.2.1 and Theorems8.5.3.2--8.5.3.3 describe the LMR
dual-degeneracy equations and determinant component. Sections8.8--8.10
separate representation-theoretic approaches and occurrence limitations.
The book is cited, not redistributed in this package.

\paragraph{Beauville and determinantal surfaces.}
Beauville, \emph{Determinantal hypersurfaces}, Michigan Mathematical
Journal48 (2000),39--64, Theorem A, Proposition1.11 and Corollary1.12,
identify the ACM cokernel and its linear resolution; the converse in
Corollary1.12 concerns a smooth hypersurface and its line bundle.
These are the precise sheaf-theoretic inputs behind the proposed
Noether--Lefschetz route. Leal--Lozano Huerta--Vite,
\emph{The Noether--Lefschetz locus of surfaces in $\mathbb P^3$
formed by determinantal surfaces},
\href{https://arxiv.org/abs/2303.09028}{arXiv:2303.09028v3},
Corollary3.1 gives the determinantal-locus dimension, and Theorem2/Table2
give degree320112 for the determinantal-quartic divisor. The latter
corrects the older doubled value640224; it is not a new computation in
this run and supplies no low-degree nonic defining equation.

\paragraph{Recent product-plus-power comparison.}
Dutta--Gesmundo--Ikenmeyer--Jindal--Lysikov,
\emph{Geometric complexity theory for product-plus-power}, Journal of
Symbolic Computation 132 (2026), 102458,
\href{https://doi.org/10.1016/j.jsc.2025.102458}{DOI:10.1016/j.jsc.2025.102458},
Theorem 1.4 and its detailed Theorem 4.10, give a multiplicity comparison
$\le4<5$ between the product-plus-power and power-sum orbit closures.
The family is $\lambda=(5d-1,1)+((d+1)\times(10d))$ for $d\ge3$.
Proposition 4.4 computes the relevant orbit representation structure;
the transport and closure argument is specific to those polynomials.
This recent precedent does not compute either padded-permanent versus
determinant restriction image in the present problem. The article was
available online on 28 May 2025 and is assigned to the 2026 volume.

\paragraph{Low-degree Chow equations.}
Cheung--Ikenmeyer--Mkrtchyan,
\href{https://arxiv.org/abs/1509.03944}{arXiv:1509.03944},
Theorem6(a) proves injectivity of the $(5,6)$ map; Theorem3 is the
McKay propagation result and Remark7 identifies the Chow ideal.
Together with lower-degree propagation these exclude Chow equations
through degree5 for$n\ge6$. Since diagonal determinants contain Chow,
this is background for starting a nonic determinant-equation search
in degree6. It is not a determinant-ideal calculation in degree6.

\paragraph{The common-linear-factor injectivity method.}
Ikenmeyer, \emph{On McKay's propagation theorem for the Foulkes
conjecture}, \href{https://arxiv.org/pdf/1509.04957}{arXiv:1509.04957v1},
Section3, equation(3.2), and Section5, Claim5.1 with its proof on
printed pp.6--7, give the common-factor map and its injectivity in
the stated range $d<n$. The whole primary PDF was read here.
QC-011 gives the elementary directional-derivative proof, including
$d=n$, and applies it to the actual ideal of forms with one linear
factor. The Fischer-inner-product lemma also appears explicitly as
Lemma2 in Section3 of the cited September2026 OpenAI paper
(source label \texttt{lem:shift}). Only this elementary lemma is
used in QC-011. The method is attributed to its predecessors.

\paragraph{Border apolarity, substitution, and viable flags.}
Buczy\'nska--Buczy\'nski, \emph{Apolarity, border rank, and
multigraded Hilbert scheme}, Duke Mathematical Journal170 (2021),
3659--3702, \href{https://arxiv.org/abs/1910.01944}{arXiv:1910.01944},
Theorem3.15 (Border Apolarity) and Theorem4.3 (Fixed Ideal Theorem),
require the relevant flat-limit component; merely passing necessary
ideal conditions does not establish smoothability.
Landsberg--Micha\l ek,
\href{https://arxiv.org/abs/1608.07486}{arXiv:1608.07486},
Proposition2.3, is the substitution reference retained in BR\_PENCIL R-01.
Conner--Harper--Landsberg, \emph{New lower bounds for matrix
multiplication and $\det_3$},
\href{https://arxiv.org/abs/1911.07981v1}{arXiv:1911.07981v1},
Theorem1.2 gives border rank17 for the polynomial coefficient tensor
of $\det_3$; sections2.2--2.5 and Proposition3.1 give the Grassmannian,
weight-space and wedge-test framework. The theorem number is pinned
to this version.

Conner--Huang--Landsberg, \emph{Bad and good news for Strassen's
laser method: Border rank of $\mathrm{perm}_3$ and strict
submultiplicativity},
\href{https://arxiv.org/abs/2009.11391v1}{arXiv:2009.11391v1},
Theorem1.1 gives the permanent coefficient-tensor border rank16.
Proposition2.1 supplies the prolongation test, Lemma2.2 the fixed-ideal
lemma, and Propositions2.3--2.4 the flag and tautological-extension
conditions. Remark2.5 distinguishes these necessary conditions from
smoothability. These principles are attributed to the source; the
new finite exclusions are the calculations in BR-004--BR-006.
The exact seventeen-term determinant curve used in BR-007 is in
section8, printed p.23, of the
\href{https://people.tamu.edu/~jml/CHLlaser5-6-21.pdf}{6 May2021 author version}.
Every coefficient of that base curve was independently checked here.
The numerical42-term expression in the following section is different
and is not an input to BR-007.

\paragraph{Post-2015 bounds and different tensor conventions.}
Farnsworth, \emph{Koszul--Young flattenings and symmetric border rank
of the determinant},
\href{https://arxiv.org/abs/1505.05079}{arXiv:1505.05079},
Theorem1.6 gives a symmetric-border-rank lower bound38. Its theorem
statement alone does not imply the same CP-border-rank lower bound.
Han--Ju--Kim, \emph{The Border Rank of the $4\times4$ Determinant
Tensor is Twelve},
\href{https://arxiv.org/abs/2510.11051v2}{arXiv:2510.11051v2},
Theorem1.1 and equation(2.1), concern the alternating tensor in
$(\mathbb C^4)^{\otimes4}$. The polynomial coefficient tensor in
$(\mathbb C^{16})^{\otimes4}$ is its Kronecker square. These bounds
therefore concern different ambient tensors.

Dole\v{z}\'alek--Micha\l ek, \emph{Nonlinear methods for tensors:
determinantal equations for secant varieties beyond cactus},
\href{https://arxiv.org/abs/2602.12762v1}{arXiv:2602.12762v1},
Theorem3.3 and Corollary3.5 give Kronecker--Koszul bounds; the
construction preceding Corollary3.9 is the quadratic tangency
flattening. Theorem1.1 separates certain secant and cactus varieties.
No bound for our tensors is inferred from this source; it records
a possible additional test. Ma\'ndziuk,
\emph{Border rank lower bounds beyond weak border apolarity},
\href{https://arxiv.org/abs/2609.12121}{arXiv:2609.12121}, was located,
but its full text could not be inspected. No unread theorem is used.

\paragraph{Recent canonical Foulkes--Howe stabilization.}
OpenAI, \emph{Quadratic stabilization of the canonical Foulkes--Howe
map}, 25 September2026, Theorem1, states surjectivity of the canonical
map $\operatorname{Sym}^b(\operatorname{Sym}^a V)\to
\operatorname{Sym}^a(\operatorname{Sym}^b V)$ for $b\ge a(a-1)$.
The primary proof in sections2--5 of
\href{https://github.com/openai/math/blob/main/preprints/Quadratic-Stabilization-of-the-Canonical-Foulkes-Howe-Map-September-25-2026/build/paper.tex}{the posted source}
was read in this run. Its dual Chow injectivity criterion requires
$n\ge30$ when the coordinate degree is6, so it does not decide
the $n=9$ or $n=23$ comparisons here. The same-date preprint
\emph{Foulkes' conjecture for the sixth symmetric power} has now been
audited in full, including its seven mathematical source sections and
both certificate programs. Theorem1.1 proves an equivariant embedding
by comparing the multiplicities in the full target; it does not assert
that the canonical map is injective at inner degrees7,8 or9. The proof
of Corollary2.4 (printed p.5, \texttt{reduction.tex}, lines189--216)
does identify the complete target weight space with the canonical
image when its first part is at least$5b$. For $b=9$ this gives the
partial threshold45, already covered by the coefficient projections
used here. It supplies no complete degree-six determinant-ideal
vanishing at size9. The posted \texttt{--check} verification of reference
hashes, tables and arithmetic bounds passed; the two C++ certificate
programs were read but not executed in this run. The separate audit
file records the corresponding base-interval and canonical-image
distinctions. No abstract embedding is substituted for a restriction
rank in the present proofs.

\paragraph{Priority scope.}
No inspected primary source was found to give the same arithmetic
seven-plane result or the same coefficient-tensor lower bounds proved
in this run. This is a negative finding within the documented search,
not proof of historical priority. The audit files give the version,
inspection scope and exact locators separately. No exhaustive novelty
claim is made for either the new calculations or the inherited runs.
